\chapter{Write and Read}
% full version: chapters/09_acet.tex

Memory in the brain has an awkward feature: memories are stored in the same connections that do the work. There is no separate hard drive. Every time the network recalls, it uses its connections, and every time it learns, it changes them. How do you write the new without spoiling the old? A computer has a separate signal for this: ``write enable.'' In the brain its role seems to be played by the \nt{ACET} radio station: acetylcholine.

\section*{Two Modes}

Imagine a network that stores several pictures: show it half of one, and it fills in the rest. That is reading. Now show it a new picture that resembles an old one. If the network fills it in out of habit, it will see the old rather than the new, and learn an error. To write the new, the network must stop trusting its memory for a while and look at the world.

\pencilfig{9}{\pencilart{9}}{One wire switches the network: write the new, or recall the old.}

Acetylcholine does exactly this, and all three of its effects are measured: it weakens the network's internal connections, through which it ``recalls,'' strengthens the signals coming from the sense organs, and makes the connections easier to change. Plenty of acetylcholine means write mode: the network listens to the world and learns. Little means read mode: the network relies on its memory and fills in.

In our model there is no level at which both go well. For writing, the internal connections must be muted; for reading, they are needed at full strength. So a switch is necessary.

\section*{How Much to Trust Your Eyes}

Engineers solved a similar problem long ago. When a ship holds its course, the navigator has a calculation of where the ship should be, and has new but imprecise observations. How much to trust each? The optimal filter answers: trust the observations the more, the less sure you are of the calculation. And the same recipe contains a second point: the less sure you are, the faster you must relearn. One knob controls both the weight of the input and the learning rate. That is exactly what acetylcholine does. Hence the hypothesis: it tells the network how uncertain its own model of the world is.

\pencilfig{124}{%
% optimal blending: the less sure the model, the more weight on the observation (and the faster it relearns)
\foreach \dx/\wm/\we/\ttop/\bbot in {0/2.4pt/0.6pt/{model is sure}/{little acetylcholine}, 4.3/0.6pt/2.4pt/{model is unsure}/{much acetylcholine}}{
  \begin{scope}[xshift=\dx cm]
  \draw[penlite=pcViolet, wash=pcViolet, rounded corners=3pt] (0,0.95) rectangle (1.3,1.45); \node[lbl, text=black] at (0.65,1.2) {memory};
  \draw[penlite=pcBlue, wash=pcBlue, rounded corners=3pt] (0,-0.25) rectangle (1.3,0.25); \node[lbl, text=black] at (0.65,0) {eyes};
  \draw[dotpen=pcAmber, wash=pcAmber] (2.95,0.6) circle (0.55); \node[lbl, text=black] at (2.95,0.6) {estimate};
  \draw[pcViolet, line width=\wm, line cap=round, -{Stealth[length=5pt]}] (1.35,1.2) -- (2.4,0.8);
  \draw[pcBlue, line width=\we, line cap=round, -{Stealth[length=5pt]}] (1.35,0) -- (2.4,0.4);
  \node[font=\small\itshape, text=pcGray, anchor=south] at (1.55,1.6) {\ttop};
  \node[lbl, anchor=north] at (1.55,-0.4) {\bbot};
  \end{scope}}
}{How much to trust your eyes. If the network is sure of its model of the world, the estimate comes almost entirely from memory. If it is unsure, from the observations, and then the network also relearns faster. According to the hypothesis, acetylcholine turns this knob.}

\section*{Division of Labor}

Recall the last chapter. When the world suddenly changes, noradrenaline, according to the hypothesis, notices the surprise. And acetylcholine keeps the network in write mode until the new situation is learned. One reports ``something has changed,'' the other ``write.''

\section*{Rewriting at Night}

During deep sleep the acetylcholine level is low. The network is cut off from the world and is in read mode. It is exactly then that the day's events are ``replayed'' in it; this is measured. That what was learned during the day is being rewritten into long-term memory at this time is a plausible hypothesis, which we will return to in the chapter on sleep.

\fullversion{9}
